% Section 3: the algebra of false agreement. The identity and bounds are
% standard facts about binary variables; their role here is to identify what
% must be measured for probe-based oversight.
\section{What Agreement Can Catch}
\label{sec:theory}

Agreement is informative only when the two readers do not fail together. This
section makes that condition precise and separates the effects of probe accuracy
and shared error.

\subsection{Setup}
\label{sec:theory:setup}

Consider only claims whose ground truth is \textsc{false}. Let $E_i=1$ when
probe $i$ incorrectly reads such a claim as \textsc{true}, and let
$p_i=\Prob(E_i=1)$ be its false-positive rate. The false-agreement rate is simply
\begin{equation}
  \FA=\Prob(E_1=1,E_2=1).
  \label{eq:fa}
\end{equation}
This is the blind spot of disagreement-based routing. Both probes make the same
error, so neither raises an objection.

\subsection{Accuracy and error correlation}
\label{sec:theory:identity}

Let $\rho$ be the Pearson correlation between the two binary error indicators.
The false-agreement rate decomposes exactly as
\begin{equation}
  \FA
  =p_1p_2
  +\rho\sqrt{p_1(1-p_1)}\sqrt{p_2(1-p_2)}.
  \label{eq:identity}
\end{equation}
This follows directly from
$\Cov(E_1,E_2)=\mathbb{E}[E_1E_2]-\mathbb{E}[E_1]\mathbb{E}[E_2]$.
The first term is the overlap expected if the probes' errors were independent;
the second is the additional overlap due to correlated error. Lower individual
error rates reduce the first term, while greater error diversity reduces the
second. Selecting a judge on accuracy alone therefore leaves out one of the two
quantities that determine the blind spot.

The two terms can also move in opposite directions. A more diverse judge may
have less-correlated errors but a higher false-positive rate, reducing the
correlation term while increasing the accuracy term. False agreement can
therefore remain unchanged even when the models' errors become less correlated.

\subsection{Bounds on false agreement}
\label{sec:theory:bounds}

The marginal error rates also constrain how much the two error sets can overlap.
The Fr\'echet bounds give
\begin{equation}
  \max(0,p_1+p_2-1)
  \leq \FA
  \leq \min(p_1,p_2).
  \label{eq:frechet}
\end{equation}
The upper bound is the error rate of the more accurate probe. It is reached when
all of that probe's errors are shared by the other. The lower bound is the
overlap forced by the two marginal error rates; it is usually zero in our
setting. Independence, $\FA=p_1p_2$, is a reference point inside this interval,
not a lower bound.

For example, if the two false-positive rates are $10\%$ and $5\%$, independence
gives $0.5\%$ false agreement, while maximal overlap gives $5\%$. The same
individual accuracies can therefore support a tenfold difference in the blind
spot, depending only on how the errors overlap.

Because the attainable correlation depends on $p_1$ and $p_2$, raw correlation
values are not always comparable across pairs with different error rates. For
positive correlations, we therefore report
\begin{equation}
  \tilde{\rho}=\frac{\rho}{\rho_{\max}},
  \qquad
  \rho_{\max}
  =\frac{\min(p_1,p_2)-p_1p_2}
  {\sqrt{p_1(1-p_1)p_2(1-p_2)}},
  \label{eq:rhonorm}
\end{equation}
which measures how much of the largest overlap permitted by the marginals is
realized; negative correlations are normalized analogously by their feasible
minimum. We report this value alongside false agreement, its independence
reference, and its Fr\'echet upper bound. These are standard consequences of
binary covariance and the Fr\'echet bounds \citep{frechet1951tableaux}; our
contribution is to use them to define and measure the blind spot of cross-model
internal readouts.

\subsection{Coverage from disagreement}
\label{sec:theory:coverage}

The same quantities describe what a disagreement router catches. Let
$p_{\mathrm{strong}}=\min(p_1,p_2)$ be the false-positive rate of the more
accurate probe. Of that probe's errors, the fraction exposed because the other
probe disagrees is
\begin{equation}
  C=1-\frac{\FA}{p_{\mathrm{strong}}}.
  \label{eq:coverage}
\end{equation}
At maximal overlap, $C=0$, so the second reader adds nothing. As overlap falls,
coverage rises. Disagreement identifies an item requiring review, however, not
which reader is correct. It therefore buys a routed check rather than an
automatic correction.
